% ECONOMICS_PROBLEM: {"id": "O38-542", "title": "Democratizing innovation", "jel_code": "O38", "source_id": "OP-0412"}
% !TeX program = xelatex
\documentclass[11pt,a4paper]{article}
\usepackage[margin=23mm]{geometry}
\usepackage{fontspec}
\setmainfont{Latin Modern Roman}
\usepackage{amsmath,amssymb,mathtools}
\usepackage{xcolor,enumitem,booktabs,longtable,array,xurl,hyperref}
\definecolor{muted}{HTML}{53636C}
\hypersetup{colorlinks=true,urlcolor=blue,linkcolor=blue}
\setlength{\parindent}{0pt}
\setlength{\parskip}{5pt}
\newcommand{\R}{\mathbb R}
\newcommand{\E}{\mathbb E}
\newcommand{\N}{\mathbb N}
\newcommand{\ind}{\mathbf 1}
\newcommand{\Var}{\operatorname{Var}}
\newcommand{\Cov}{\operatorname{Cov}}
\newcommand{\supp}{\operatorname{supp}}
\DeclareMathOperator*{\argmin}{arg\,min}
\DeclareMathOperator*{\argmax}{arg\,max}
\newcommand{\entrymeta}[3]{{\sffamily\small\color{muted}\textbf{\texttt{#1}}\quad|\quad #2\par #3\par}}
\newcommand{\Problemref}[1]{\texttt{#1}}
\newcommand{\JELref}[2]{\texttt{#2} (JEL \texttt{#1})}
\begin{document}
\section*{Democratizing innovation}
\textbf{Problem ID:} \texttt{O38-542}\quad \textbf{JEL:} \texttt{O38}\par
% BEGIN ECONOMICS STATEMENT
\label{problem:OP-0412}
{\sffamily\small\color{muted}Original subject group: Innovation and AI economics\par}
\label{definitions:problem:30007694}
\textbf{Audit classification:} Explicit published open question. The publication explicitly asks about access democratization. The original term change in the gap had a sign ambiguity and access was not distinguished from use.
\subsection*{Definitions and precise research target}
Consider firms eligible for a public AI research-access program, partitioned before treatment into small firms \(G=s\) and resource-rich incumbents \(G=l\) using a stated employment threshold. Let \(A=1\) provide an identical subsidized model, computation allowance and training package, while \(A=0\) retains existing access. Over three years, define \(Y_i(a)\) as the value of independently validated new products from firm \(i\), with zero for nonparticipation or unsuccessful research. Group effects are
\[
\tau_g=E[Y_i(1)-Y_i(0)\mid G=g],\qquad g\in\{s,l\},
\]
and the closure of the incumbent-minus-small average innovation gap is \(\Gamma=\tau_s-\tau_l\). This gap does not itself measure aggregate welfare, so separately record entry, exit, innovation expenditure and product diffusion. Use program eligibility lotteries or clustered rollout with recorded access compliance; account for knowledge spillovers and changes in incumbent competition within product markets. Define a market-level policy effect by assigning access to all eligible firms in a market and measuring total validated value net of program resource costs. Which access and support arrangements make \(\Gamma>0\) while increasing this aggregate outcome? The scope is eligible firms in the selected sectors, with uncertainty about commercialization and crowding out explicitly measured rather than an unrestricted claim that AI democratizes innovation.

\textbf{Definitions, assumptions and mathematical qualifications.} Group membership is fixed at baseline. Define \(\mu_g(a)=E[Y_i(a)\mid G=g]\). The incumbent-minus-small mean gap is \(D(a)=\mu_l(a)-\mu_s(a)\), and \(\Gamma=D(0)-D(1)=\tau_s-\tau_l\) is gap closure; it is not the signed change \(D(1)-D(0)\). Assignment to the access package, actual use and successful innovation are separate variables. Define the policy effect at the product-market level, with market-wide eligibility and potential knowledge spillovers, to evaluate commercialization and displacement. A larger small-firm effect alone need not imply a positive aggregate resource-adjusted effect.
\subsection*{Variants and assumption changes}
\begin{enumerate}
\item \textbf{Editorial, intention to treat.} Lottery assignment identifies group-specific package effects including nonuse, under groupwise assignment positivity; estimate gap closure without an exclusion assumption.
\item \textbf{Editorial, policy saturation.} Randomize whole markets to universal eligible-firm access and compare total validated value less program and research costs, allowing incumbents to respond and products to displace each other.
\end{enumerate}
\subsection*{References and source audit}
\textbf{Checked source.} 2024 Stanford draft, Section 3.2, PDF pp. 10--11. Full primary draft and cover retrieved. \url{https://digitaleconomy.stanford.edu/app/uploads/2024/11/ETAI-White-Paper.pdf}.
\begin{enumerate}\raggedright
\item\hypertarget{definitions:source:30007694:1}{} Erik Brynjolfsson; Anton Korinek; Ajay Agrawal. 2024. \emph{The Economics of Transformative AI: A Research Agenda}. 2024 Stanford draft, Section 3.2, PDF pp. 10--11.
\url{https://digitaleconomy.stanford.edu/app/uploads/2024/11/ETAI-White-Paper.pdf}.
\end{enumerate}

\subsection*{Common conventions: Source mathematical conventions}

These conventions apply to each entry unless explicitly replaced.

\textbf{Models and identification.} A primitive model \(m\) specifies all probability laws, feasible choices, information, equilibrium and measurement rules used by its target. The admissible class \(\mathcal M\) is fixed independently of outcomes. If \(P_O(m)\) is its observable probability law and \(\tau(m)\) the target, its identified set at observable law \(P\) is
\[
 \mathcal I_\tau(P)=\{\tau(m):m\in\mathcal M,\ P_O(m)=P\}.
\]
Point identification means that this set is a singleton. An empty set signals incompatibility of data and assumptions. Coordinate bounds need not describe the joint set. Bounds are sharp only if every retained target is attainable or arbitrarily approachable by compatible models. Finite bounds require explicit restrictions on unobserved quantities; unrestricted confounding, hidden wealth, extrapolation or arbitrary equilibrium selection can yield uninformative sets.

\textbf{Causal interventions.} A potential outcome \(Y_i(p)\) is the outcome under a complete policy regime \(p\), including timing, financing, information and market spillovers. The target population and units are fixed before treatment. With interference, use the complete assignment vector or a justified exposure mapping; individual no-interference notation alone cannot identify an economy-wide intervention. A difference of mean potential outcomes does not require the cross-world joint distribution. Conditioning on post-treatment migration, survival, enrollment or employment changes the estimand and needs separate justification.

\textbf{Statistical statements.} An estimator, confidence set, forecast or test specifies a sampling law, model class and loss/coverage criterion. Uniform \((1-\alpha)\)-coverage means \(\inf_{m\in\mathcal M}\Pr_m\{\tau(m)\in C_n\}\geq1-\alpha\), or an explicitly asymptotic version. The whole parameter space trivially has coverage, so informativeness requires a stated width, risk or power objective. A forecasting improvement is relative to a named benchmark, information set and evaluation population.

\textbf{Equilibrium and optimization.} An equilibrium comprises feasible plans, optimality, beliefs where needed, budget constraints and all market/resource clearing. The primitives or a stated benchmark define each correspondence \(\mathcal E(p;m)\). Nonemptiness is assumed or proved, never inferred just from compact policy sets. A selected-equilibrium target uses a declared selection rule; a robust target quantifies over all permitted equilibria. Use suprema or infima unless attainment is established. An \(\varepsilon\)-optimal policy is feasible and within \(\varepsilon\) of the finite optimum value. Report an empty feasible set separately.

\textbf{Welfare and accounting.} Normative utility, interpersonal normalization, social weights, numeraire, discounting and the use of public receipts are primitives. Behavioral utility need not coincide with the welfare criterion. Household welfare includes ownership income and rents once; adding profits again to compensating variation generally double-counts them. A monetary equivalent is defined by an explicit transfer and price convention, with monotonicity and range conditions. Average effects, mechanism contributions, transfers and welfare losses are different targets. Infinite sums require convergence and dynamic feasible plans require terminal or no-Ponzi conditions.

\textbf{Source versions.} A working-paper date, revision date and journal publication year are distinguished. A DOI for a journal version is not automatically the DOI of an earlier manuscript. Locators refer to the stated version and printed pages where specified. A metadata-only check does not verify a claim about a full-text conclusion. Every entry's source subsection narrows the provenance claim accordingly.
% END ECONOMICS STATEMENT
\end{document}
